% Options for packages loaded elsewhere
\PassOptionsToPackage{unicode}{hyperref}
\PassOptionsToPackage{hyphens}{url}
\documentclass[
]{article}
\usepackage{xcolor}
\usepackage[margin=1in]{geometry}
\usepackage{amsmath,amssymb}
\setcounter{secnumdepth}{5}
\usepackage{iftex}
\ifPDFTeX
  \usepackage[T1]{fontenc}
  \usepackage[utf8]{inputenc}
  \usepackage{textcomp} % provide euro and other symbols
\else % if luatex or xetex
  \usepackage{unicode-math} % this also loads fontspec
  \defaultfontfeatures{Scale=MatchLowercase}
  \defaultfontfeatures[\rmfamily]{Ligatures=TeX,Scale=1}
\fi
\usepackage{lmodern}
\ifPDFTeX\else
  % xetex/luatex font selection
\fi
% Use upquote if available, for straight quotes in verbatim environments
\IfFileExists{upquote.sty}{\usepackage{upquote}}{}
\IfFileExists{microtype.sty}{% use microtype if available
  \usepackage[]{microtype}
  \UseMicrotypeSet[protrusion]{basicmath} % disable protrusion for tt fonts
}{}
\makeatletter
\@ifundefined{KOMAClassName}{% if non-KOMA class
  \IfFileExists{parskip.sty}{%
    \usepackage{parskip}
  }{% else
    \setlength{\parindent}{0pt}
    \setlength{\parskip}{6pt plus 2pt minus 1pt}}
}{% if KOMA class
  \KOMAoptions{parskip=half}}
\makeatother
\usepackage{color}
\usepackage{fancyvrb}
\newcommand{\VerbBar}{|}
\newcommand{\VERB}{\Verb[commandchars=\\\{\}]}
\DefineVerbatimEnvironment{Highlighting}{Verbatim}{commandchars=\\\{\}}
% Add ',fontsize=\small' for more characters per line
\usepackage{framed}
\definecolor{shadecolor}{RGB}{248,248,248}
\newenvironment{Shaded}{\begin{snugshade}}{\end{snugshade}}
\newcommand{\AlertTok}[1]{\textcolor[rgb]{0.94,0.16,0.16}{#1}}
\newcommand{\AnnotationTok}[1]{\textcolor[rgb]{0.56,0.35,0.01}{\textbf{\textit{#1}}}}
\newcommand{\AttributeTok}[1]{\textcolor[rgb]{0.13,0.29,0.53}{#1}}
\newcommand{\BaseNTok}[1]{\textcolor[rgb]{0.00,0.00,0.81}{#1}}
\newcommand{\BuiltInTok}[1]{#1}
\newcommand{\CharTok}[1]{\textcolor[rgb]{0.31,0.60,0.02}{#1}}
\newcommand{\CommentTok}[1]{\textcolor[rgb]{0.56,0.35,0.01}{\textit{#1}}}
\newcommand{\CommentVarTok}[1]{\textcolor[rgb]{0.56,0.35,0.01}{\textbf{\textit{#1}}}}
\newcommand{\ConstantTok}[1]{\textcolor[rgb]{0.56,0.35,0.01}{#1}}
\newcommand{\ControlFlowTok}[1]{\textcolor[rgb]{0.13,0.29,0.53}{\textbf{#1}}}
\newcommand{\DataTypeTok}[1]{\textcolor[rgb]{0.13,0.29,0.53}{#1}}
\newcommand{\DecValTok}[1]{\textcolor[rgb]{0.00,0.00,0.81}{#1}}
\newcommand{\DocumentationTok}[1]{\textcolor[rgb]{0.56,0.35,0.01}{\textbf{\textit{#1}}}}
\newcommand{\ErrorTok}[1]{\textcolor[rgb]{0.64,0.00,0.00}{\textbf{#1}}}
\newcommand{\ExtensionTok}[1]{#1}
\newcommand{\FloatTok}[1]{\textcolor[rgb]{0.00,0.00,0.81}{#1}}
\newcommand{\FunctionTok}[1]{\textcolor[rgb]{0.13,0.29,0.53}{\textbf{#1}}}
\newcommand{\ImportTok}[1]{#1}
\newcommand{\InformationTok}[1]{\textcolor[rgb]{0.56,0.35,0.01}{\textbf{\textit{#1}}}}
\newcommand{\KeywordTok}[1]{\textcolor[rgb]{0.13,0.29,0.53}{\textbf{#1}}}
\newcommand{\NormalTok}[1]{#1}
\newcommand{\OperatorTok}[1]{\textcolor[rgb]{0.81,0.36,0.00}{\textbf{#1}}}
\newcommand{\OtherTok}[1]{\textcolor[rgb]{0.56,0.35,0.01}{#1}}
\newcommand{\PreprocessorTok}[1]{\textcolor[rgb]{0.56,0.35,0.01}{\textit{#1}}}
\newcommand{\RegionMarkerTok}[1]{#1}
\newcommand{\SpecialCharTok}[1]{\textcolor[rgb]{0.81,0.36,0.00}{\textbf{#1}}}
\newcommand{\SpecialStringTok}[1]{\textcolor[rgb]{0.31,0.60,0.02}{#1}}
\newcommand{\StringTok}[1]{\textcolor[rgb]{0.31,0.60,0.02}{#1}}
\newcommand{\VariableTok}[1]{\textcolor[rgb]{0.00,0.00,0.00}{#1}}
\newcommand{\VerbatimStringTok}[1]{\textcolor[rgb]{0.31,0.60,0.02}{#1}}
\newcommand{\WarningTok}[1]{\textcolor[rgb]{0.56,0.35,0.01}{\textbf{\textit{#1}}}}
\usepackage{graphicx}
\makeatletter
\newsavebox\pandoc@box
\newcommand*\pandocbounded[1]{% scales image to fit in text height/width
  \sbox\pandoc@box{#1}%
  \Gscale@div\@tempa{\textheight}{\dimexpr\ht\pandoc@box+\dp\pandoc@box\relax}%
  \Gscale@div\@tempb{\linewidth}{\wd\pandoc@box}%
  \ifdim\@tempb\p@<\@tempa\p@\let\@tempa\@tempb\fi% select the smaller of both
  \ifdim\@tempa\p@<\p@\scalebox{\@tempa}{\usebox\pandoc@box}%
  \else\usebox{\pandoc@box}%
  \fi%
}
% Set default figure placement to htbp
\def\fps@figure{htbp}
\makeatother
\setlength{\emergencystretch}{3em} % prevent overfull lines
\providecommand{\tightlist}{%
  \setlength{\itemsep}{0pt}\setlength{\parskip}{0pt}}
\usepackage{bookmark}
\IfFileExists{xurl.sty}{\usepackage{xurl}}{} % add URL line breaks if available
\urlstyle{same}
\hypersetup{
  pdftitle={Task 3: Interpretability - Linear Models},
  pdfauthor={Mario Álvarez Martínez, Marcos Carrasco Panadero, Sergio Samaniego Hernández},
  hidelinks,
  pdfcreator={LaTeX via pandoc}}

\title{Task 3: Interpretability - Linear Models}
\usepackage{etoolbox}
\makeatletter
\providecommand{\subtitle}[1]{% add subtitle to \maketitle
  \apptocmd{\@title}{\par {\large #1 \par}}{}{}
}
\makeatother
\subtitle{Reliable Models - EDM-GCD}
\author{Mario Álvarez Martínez, Marcos Carrasco Panadero, Sergio
Samaniego Hernández}
\date{2026-04-22}

\begin{document}
\maketitle

{
\setcounter{tocdepth}{3}
\tableofcontents
}
\section{Introduction}\label{introduction}

\subsection{Context and Motivation}\label{context-and-motivation}

In this task, we will address the topic of \textbf{Interpretability} of
machine learning models. Interpretability is fundamental in many fields
such as medicine, finance, public administration, etc., where not only
does the model's predictive accuracy matter, but also its ability to
explain \textbf{why} it makes certain decisions.

\textbf{Key question}: How can we understand which features influence
our model's predictions?

\subsubsection{Objectives of this
practice}\label{objectives-of-this-practice}

\begin{itemize}
\tightlist
\item
  Understand the principles of \textbf{interpretability} in linear
  regression models
\item
  Learn to \textbf{interpret model coefficients}
\item
  Evaluate the \textbf{validity of interpretations} through diagnostics
\item
  Visualize and communicate results clearly
\item
  Understand the \textbf{limitations} of linear models
\end{itemize}

\subsection{The Problem: Predicting Bicycle
Rentals}\label{the-problem-predicting-bicycle-rentals}

We will use the \textbf{Bike-Sharing} dataset from the Capital Bikeshare
system in Washington D.C. (2011-2012).

\subsubsection{Research question}\label{research-question}

\emph{What are the factors that influence the number of bicycles rented
daily and what is the impact of each of these factors?}

\subsubsection{Dataset variables}\label{dataset-variables}

The dataset contains aggregated \textbf{daily} information about bicycle
rentals:

\textbf{Dependent variable (target):} - \texttt{cnt}: Total number of
bicycles rented (casual + registered)

\textbf{Independent variables (features):} - \texttt{season}: Season (1:
spring, 2: summer, 3: fall, 4: winter) - \texttt{yr}: Year (0: 2011, 1:
2012) - \texttt{mnth}: Month (1-12) - \texttt{holiday}: Is holiday (0/1)
- \texttt{weekday}: Day of week (0-6) - \texttt{workingday}: Is working
day (0/1) - \texttt{weathersit}: Weather condition (1: clear, 2: cloudy,
3: light rain, 4: heavy rain) - \texttt{temp}: Normalized temperature
(0-1) - \texttt{atemp}: Normalized apparent temperature (0-1) -
\texttt{hum}: Normalized humidity (0-1) - \texttt{windspeed}: Normalized
wind speed (0-1)

\subsection{Initial hypotheses}\label{initial-hypotheses}

Based on theory and common sense, we expect:

\begin{enumerate}
\def\labelenumi{\arabic{enumi}.}
\tightlist
\item
  \textbf{Positive effect of temperature}: Warmer days → more rentals
\item
  \textbf{Negative effect of bad weather}: Rain/snow → fewer rentals
\item
  \textbf{Effect of day of week}: Differences between weekdays and
  weekends
\item
  \textbf{Effect of season}: Variations according to time of year
\item
  \textbf{Effect of year}: Possible growth from 2011 to 2012
\end{enumerate}

\begin{center}\rule{0.5\linewidth}{0.5pt}\end{center}

\section{The Linear Model}\label{the-linear-model}

\subsection{Loading and exploring the
dataset}\label{loading-and-exploring-the-dataset}

\begin{Shaded}
\begin{Highlighting}[]
\CommentTok{\# Load the dataset}
\NormalTok{data }\OtherTok{\textless{}{-}} \FunctionTok{read.csv}\NormalTok{(}\StringTok{"day.csv"}\NormalTok{)}

\CommentTok{\# Initial exploration}
\FunctionTok{head}\NormalTok{(data, }\DecValTok{10}\NormalTok{)}
\end{Highlighting}
\end{Shaded}

\begin{verbatim}
##    instant     dteday season yr mnth holiday weekday workingday weathersit
## 1        1 2011-01-01      1  0    1       0       6          0          2
## 2        2 2011-01-02      1  0    1       0       0          0          2
## 3        3 2011-01-03      1  0    1       0       1          1          1
## 4        4 2011-01-04      1  0    1       0       2          1          1
## 5        5 2011-01-05      1  0    1       0       3          1          1
## 6        6 2011-01-06      1  0    1       0       4          1          1
## 7        7 2011-01-07      1  0    1       0       5          1          2
## 8        8 2011-01-08      1  0    1       0       6          0          2
## 9        9 2011-01-09      1  0    1       0       0          0          1
## 10      10 2011-01-10      1  0    1       0       1          1          1
##        temp    atemp      hum windspeed casual registered  cnt
## 1  0.344167 0.363625 0.805833 0.1604460    331        654  985
## 2  0.363478 0.353739 0.696087 0.2485390    131        670  801
## 3  0.196364 0.189405 0.437273 0.2483090    120       1229 1349
## 4  0.200000 0.212122 0.590435 0.1602960    108       1454 1562
## 5  0.226957 0.229270 0.436957 0.1869000     82       1518 1600
## 6  0.204348 0.233209 0.518261 0.0895652     88       1518 1606
## 7  0.196522 0.208839 0.498696 0.1687260    148       1362 1510
## 8  0.165000 0.162254 0.535833 0.2668040     68        891  959
## 9  0.138333 0.116175 0.434167 0.3619500     54        768  822
## 10 0.150833 0.150888 0.482917 0.2232670     41       1280 1321
\end{verbatim}

\begin{Shaded}
\begin{Highlighting}[]
\FunctionTok{dim}\NormalTok{(data)}
\end{Highlighting}
\end{Shaded}

\begin{verbatim}
## [1] 731  16
\end{verbatim}

\begin{Shaded}
\begin{Highlighting}[]
\FunctionTok{str}\NormalTok{(data)}
\end{Highlighting}
\end{Shaded}

\begin{verbatim}
## 'data.frame':    731 obs. of  16 variables:
##  $ instant   : int  1 2 3 4 5 6 7 8 9 10 ...
##  $ dteday    : chr  "2011-01-01" "2011-01-02" "2011-01-03" "2011-01-04" ...
##  $ season    : int  1 1 1 1 1 1 1 1 1 1 ...
##  $ yr        : int  0 0 0 0 0 0 0 0 0 0 ...
##  $ mnth      : int  1 1 1 1 1 1 1 1 1 1 ...
##  $ holiday   : int  0 0 0 0 0 0 0 0 0 0 ...
##  $ weekday   : int  6 0 1 2 3 4 5 6 0 1 ...
##  $ workingday: int  0 0 1 1 1 1 1 0 0 1 ...
##  $ weathersit: int  2 2 1 1 1 1 2 2 1 1 ...
##  $ temp      : num  0.344 0.363 0.196 0.2 0.227 ...
##  $ atemp     : num  0.364 0.354 0.189 0.212 0.229 ...
##  $ hum       : num  0.806 0.696 0.437 0.59 0.437 ...
##  $ windspeed : num  0.16 0.249 0.248 0.16 0.187 ...
##  $ casual    : int  331 131 120 108 82 88 148 68 54 41 ...
##  $ registered: int  654 670 1229 1454 1518 1518 1362 891 768 1280 ...
##  $ cnt       : int  985 801 1349 1562 1600 1606 1510 959 822 1321 ...
\end{verbatim}

\begin{Shaded}
\begin{Highlighting}[]
\CommentTok{\# Descriptive statistics of main variables}
\FunctionTok{summary}\NormalTok{(data)}
\end{Highlighting}
\end{Shaded}

\begin{verbatim}
##     instant         dteday              season            yr        
##  Min.   :  1.0   Length:731         Min.   :1.000   Min.   :0.0000  
##  1st Qu.:183.5   Class :character   1st Qu.:2.000   1st Qu.:0.0000  
##  Median :366.0   Mode  :character   Median :3.000   Median :1.0000  
##  Mean   :366.0                      Mean   :2.497   Mean   :0.5007  
##  3rd Qu.:548.5                      3rd Qu.:3.000   3rd Qu.:1.0000  
##  Max.   :731.0                      Max.   :4.000   Max.   :1.0000  
##       mnth          holiday           weekday        workingday   
##  Min.   : 1.00   Min.   :0.00000   Min.   :0.000   Min.   :0.000  
##  1st Qu.: 4.00   1st Qu.:0.00000   1st Qu.:1.000   1st Qu.:0.000  
##  Median : 7.00   Median :0.00000   Median :3.000   Median :1.000  
##  Mean   : 6.52   Mean   :0.02873   Mean   :2.997   Mean   :0.684  
##  3rd Qu.:10.00   3rd Qu.:0.00000   3rd Qu.:5.000   3rd Qu.:1.000  
##  Max.   :12.00   Max.   :1.00000   Max.   :6.000   Max.   :1.000  
##    weathersit         temp             atemp              hum        
##  Min.   :1.000   Min.   :0.05913   Min.   :0.07907   Min.   :0.0000  
##  1st Qu.:1.000   1st Qu.:0.33708   1st Qu.:0.33784   1st Qu.:0.5200  
##  Median :1.000   Median :0.49833   Median :0.48673   Median :0.6267  
##  Mean   :1.395   Mean   :0.49538   Mean   :0.47435   Mean   :0.6279  
##  3rd Qu.:2.000   3rd Qu.:0.65542   3rd Qu.:0.60860   3rd Qu.:0.7302  
##  Max.   :3.000   Max.   :0.86167   Max.   :0.84090   Max.   :0.9725  
##    windspeed           casual         registered        cnt      
##  Min.   :0.02239   Min.   :   2.0   Min.   :  20   Min.   :  22  
##  1st Qu.:0.13495   1st Qu.: 315.5   1st Qu.:2497   1st Qu.:3152  
##  Median :0.18097   Median : 713.0   Median :3662   Median :4548  
##  Mean   :0.19049   Mean   : 848.2   Mean   :3656   Mean   :4504  
##  3rd Qu.:0.23321   3rd Qu.:1096.0   3rd Qu.:4776   3rd Qu.:5956  
##  Max.   :0.50746   Max.   :3410.0   Max.   :6946   Max.   :8714
\end{verbatim}

\begin{Shaded}
\begin{Highlighting}[]
\CommentTok{\# Focus on the target variable}
\FunctionTok{cat}\NormalTok{(}\StringTok{"Statistics of cnt (daily rentals):}\SpecialCharTok{\textbackslash{}n}\StringTok{"}\NormalTok{)}
\end{Highlighting}
\end{Shaded}

\begin{verbatim}
## Statistics of cnt (daily rentals):
\end{verbatim}

\begin{Shaded}
\begin{Highlighting}[]
\FunctionTok{cat}\NormalTok{(}\StringTok{"Mean:"}\NormalTok{, }\FunctionTok{mean}\NormalTok{(data}\SpecialCharTok{$}\NormalTok{cnt), }\StringTok{"}\SpecialCharTok{\textbackslash{}n}\StringTok{"}\NormalTok{)}
\end{Highlighting}
\end{Shaded}

\begin{verbatim}
## Mean: 4504.349
\end{verbatim}

\begin{Shaded}
\begin{Highlighting}[]
\FunctionTok{cat}\NormalTok{(}\StringTok{"Standard deviation:"}\NormalTok{, }\FunctionTok{sd}\NormalTok{(data}\SpecialCharTok{$}\NormalTok{cnt), }\StringTok{"}\SpecialCharTok{\textbackslash{}n}\StringTok{"}\NormalTok{)}
\end{Highlighting}
\end{Shaded}

\begin{verbatim}
## Standard deviation: 1937.211
\end{verbatim}

\begin{Shaded}
\begin{Highlighting}[]
\FunctionTok{cat}\NormalTok{(}\StringTok{"Minimum:"}\NormalTok{, }\FunctionTok{min}\NormalTok{(data}\SpecialCharTok{$}\NormalTok{cnt), }\StringTok{"}\SpecialCharTok{\textbackslash{}n}\StringTok{"}\NormalTok{)}
\end{Highlighting}
\end{Shaded}

\begin{verbatim}
## Minimum: 22
\end{verbatim}

\begin{Shaded}
\begin{Highlighting}[]
\FunctionTok{cat}\NormalTok{(}\StringTok{"Maximum:"}\NormalTok{, }\FunctionTok{max}\NormalTok{(data}\SpecialCharTok{$}\NormalTok{cnt), }\StringTok{"}\SpecialCharTok{\textbackslash{}n}\StringTok{"}\NormalTok{)}
\end{Highlighting}
\end{Shaded}

\begin{verbatim}
## Maximum: 8714
\end{verbatim}

\subsection{Data preparation}\label{data-preparation}

\begin{Shaded}
\begin{Highlighting}[]
\CommentTok{\# IMPORTANT: Convert categorical variables to factors for better interpretation}

\CommentTok{\# Create copies to work with}
\NormalTok{data\_model }\OtherTok{\textless{}{-}}\NormalTok{ data}

\CommentTok{\# Convert categorical variables to factors}
\CommentTok{\# Convert categorical variables to factors}
\NormalTok{data\_model}\SpecialCharTok{$}\NormalTok{season }\OtherTok{\textless{}{-}} \FunctionTok{factor}\NormalTok{(data\_model}\SpecialCharTok{$}\NormalTok{season, }
                            \AttributeTok{levels =} \FunctionTok{c}\NormalTok{(}\DecValTok{1}\NormalTok{, }\DecValTok{2}\NormalTok{, }\DecValTok{3}\NormalTok{, }\DecValTok{4}\NormalTok{),}
                            \AttributeTok{labels =} \FunctionTok{c}\NormalTok{(}\StringTok{"Spring"}\NormalTok{, }\StringTok{"Summer"}\NormalTok{, }\StringTok{"Fall"}\NormalTok{, }\StringTok{"Winter"}\NormalTok{))}
\NormalTok{data\_model}\SpecialCharTok{$}\NormalTok{yr }\OtherTok{\textless{}{-}} \FunctionTok{factor}\NormalTok{(data\_model}\SpecialCharTok{$}\NormalTok{yr, }
                        \AttributeTok{levels =} \FunctionTok{c}\NormalTok{(}\DecValTok{0}\NormalTok{, }\DecValTok{1}\NormalTok{),}
                        \AttributeTok{labels =} \FunctionTok{c}\NormalTok{(}\StringTok{"2011"}\NormalTok{, }\StringTok{"2012"}\NormalTok{))}
\NormalTok{data\_model}\SpecialCharTok{$}\NormalTok{mnth }\OtherTok{\textless{}{-}} \FunctionTok{factor}\NormalTok{(data\_model}\SpecialCharTok{$}\NormalTok{mnth)}
\NormalTok{data\_model}\SpecialCharTok{$}\NormalTok{weekday }\OtherTok{\textless{}{-}} \FunctionTok{factor}\NormalTok{(data\_model}\SpecialCharTok{$}\NormalTok{weekday)}
\NormalTok{data\_model}\SpecialCharTok{$}\NormalTok{weathersit }\OtherTok{\textless{}{-}} \FunctionTok{factor}\NormalTok{(data\_model}\SpecialCharTok{$}\NormalTok{weathersit,}
                                \AttributeTok{levels =} \FunctionTok{c}\NormalTok{(}\DecValTok{1}\NormalTok{, }\DecValTok{2}\NormalTok{, }\DecValTok{3}\NormalTok{, }\DecValTok{4}\NormalTok{),}
                                \AttributeTok{labels =} \FunctionTok{c}\NormalTok{(}\StringTok{"Clear"}\NormalTok{, }\StringTok{"Mist"}\NormalTok{, }\StringTok{"LightRain"}\NormalTok{, }\StringTok{"HeavyRain"}\NormalTok{))}
\NormalTok{data\_model}\SpecialCharTok{$}\NormalTok{holiday }\OtherTok{\textless{}{-}} \FunctionTok{factor}\NormalTok{(data\_model}\SpecialCharTok{$}\NormalTok{holiday)}
\NormalTok{data\_model}\SpecialCharTok{$}\NormalTok{workingday }\OtherTok{\textless{}{-}} \FunctionTok{factor}\NormalTok{(data\_model}\SpecialCharTok{$}\NormalTok{workingday)}

\CommentTok{\# Visualize the distribution of the target variable}
\NormalTok{p1 }\OtherTok{\textless{}{-}} \FunctionTok{ggplot}\NormalTok{(data\_model, }\FunctionTok{aes}\NormalTok{(}\AttributeTok{x =}\NormalTok{ cnt)) }\SpecialCharTok{+}
  \FunctionTok{geom\_histogram}\NormalTok{(}\AttributeTok{bins =} \DecValTok{30}\NormalTok{, }\AttributeTok{fill =} \StringTok{"steelblue"}\NormalTok{, }\AttributeTok{color =} \StringTok{"black"}\NormalTok{, }\AttributeTok{alpha =} \FloatTok{0.7}\NormalTok{) }\SpecialCharTok{+}
  \FunctionTok{labs}\NormalTok{(}\AttributeTok{title =} \StringTok{"Distribution of daily rentals (cnt)"}\NormalTok{,}
       \AttributeTok{x =} \StringTok{"Number of bicycles rented"}\NormalTok{,}
       \AttributeTok{y =} \StringTok{"Frequency"}\NormalTok{) }\SpecialCharTok{+}
  \FunctionTok{theme\_minimal}\NormalTok{() }\SpecialCharTok{+}
  \FunctionTok{theme}\NormalTok{(}\AttributeTok{plot.title =} \FunctionTok{element\_text}\NormalTok{(}\AttributeTok{size =} \DecValTok{12}\NormalTok{, }\AttributeTok{face =} \StringTok{"bold"}\NormalTok{))}

\NormalTok{p2 }\OtherTok{\textless{}{-}} \FunctionTok{ggplot}\NormalTok{(data\_model, }\FunctionTok{aes}\NormalTok{(}\AttributeTok{sample =}\NormalTok{ cnt)) }\SpecialCharTok{+}
  \FunctionTok{stat\_qq}\NormalTok{(}\AttributeTok{color =} \StringTok{"steelblue"}\NormalTok{) }\SpecialCharTok{+}
  \FunctionTok{stat\_qq\_line}\NormalTok{(}\AttributeTok{color =} \StringTok{"red"}\NormalTok{) }\SpecialCharTok{+}
  \FunctionTok{labs}\NormalTok{(}\AttributeTok{title =} \StringTok{"Q{-}Q Plot: Normality of cnt"}\NormalTok{,}
       \AttributeTok{x =} \StringTok{"Theoretical Quantiles"}\NormalTok{,}
       \AttributeTok{y =} \StringTok{"Sample Quantiles"}\NormalTok{) }\SpecialCharTok{+}
  \FunctionTok{theme\_minimal}\NormalTok{() }\SpecialCharTok{+}
  \FunctionTok{theme}\NormalTok{(}\AttributeTok{plot.title =} \FunctionTok{element\_text}\NormalTok{(}\AttributeTok{size =} \DecValTok{12}\NormalTok{, }\AttributeTok{face =} \StringTok{"bold"}\NormalTok{))}

\FunctionTok{grid.arrange}\NormalTok{(p1, p2, }\AttributeTok{ncol =} \DecValTok{2}\NormalTok{)}
\end{Highlighting}
\end{Shaded}

\pandocbounded{\includegraphics[keepaspectratio]{Practica_4_Interpretability_EN_files/figure-latex/data_preparation-1.pdf}}

\begin{Shaded}
\begin{Highlighting}[]
\CommentTok{\# MEMORY }\AlertTok{NOTE}\CommentTok{:}
\CommentTok{\# We observe that the distribution of cnt is approximately normal, }
\CommentTok{\# which justifies the use of linear regression.}
\end{Highlighting}
\end{Shaded}

\subsection{Correlation analysis}\label{correlation-analysis}

\begin{Shaded}
\begin{Highlighting}[]
\CommentTok{\# Select only numeric variables for correlation matrix}
\NormalTok{numeric\_cols }\OtherTok{\textless{}{-}} \FunctionTok{c}\NormalTok{(}\StringTok{"temp"}\NormalTok{, }\StringTok{"atemp"}\NormalTok{, }\StringTok{"hum"}\NormalTok{, }\StringTok{"windspeed"}\NormalTok{, }\StringTok{"cnt"}\NormalTok{)}
\NormalTok{cor\_matrix }\OtherTok{\textless{}{-}} \FunctionTok{cor}\NormalTok{(data[, numeric\_cols])}

\CommentTok{\# Visualize correlation matrix}
\FunctionTok{corrplot}\NormalTok{(cor\_matrix, }\AttributeTok{method =} \StringTok{"circle"}\NormalTok{, }\AttributeTok{type =} \StringTok{"upper"}\NormalTok{, }
         \AttributeTok{diag =} \ConstantTok{FALSE}\NormalTok{, }\AttributeTok{addCoef.col =} \StringTok{"black"}\NormalTok{, }\AttributeTok{tl.cex =} \FloatTok{0.8}\NormalTok{,}
         \AttributeTok{title =} \StringTok{"Correlation matrix {-} Numeric variables"}\NormalTok{)}
\end{Highlighting}
\end{Shaded}

\pandocbounded{\includegraphics[keepaspectratio]{Practica_4_Interpretability_EN_files/figure-latex/correlation-1.pdf}}

\begin{Shaded}
\begin{Highlighting}[]
\CommentTok{\# MEMORY }\AlertTok{NOTE}\CommentTok{:}
\CommentTok{\# {-} Temperature (temp) and apparent temperature (atemp) are highly correlated}
\CommentTok{\# {-} There is multicollinearity between temp and atemp (r \textgreater{} 0.98)}
\CommentTok{\# {-} cnt is positively correlated with temperature}
\CommentTok{\# {-} cnt is negatively correlated with humidity and wind speed}
\end{Highlighting}
\end{Shaded}

\subsection{Fitting the Multiple Linear
Model}\label{fitting-the-multiple-linear-model}

\begin{Shaded}
\begin{Highlighting}[]
\CommentTok{\# MEMORY }\AlertTok{NOTE}\CommentTok{:}
\CommentTok{\# We fit a linear model that includes:}
\CommentTok{\# {-} Numeric variables: temperature, humidity, wind speed}
\CommentTok{\# {-} Categorical variables: season, year, weather condition}
\CommentTok{\# {-} We exclude \textquotesingle{}instant\textquotesingle{}, \textquotesingle{}dteday\textquotesingle{}, \textquotesingle{}casual\textquotesingle{} and \textquotesingle{}registered\textquotesingle{} (irrelevant/redundant variables)}
\CommentTok{\# {-} We exclude \textquotesingle{}atemp\textquotesingle{} due to multicollinearity with \textquotesingle{}temp\textquotesingle{}}

\NormalTok{model }\OtherTok{\textless{}{-}} \FunctionTok{lm}\NormalTok{(cnt }\SpecialCharTok{\textasciitilde{}}\NormalTok{ temp }\SpecialCharTok{+}\NormalTok{ hum }\SpecialCharTok{+}\NormalTok{ windspeed }\SpecialCharTok{+}\NormalTok{ season }\SpecialCharTok{+}\NormalTok{ yr }\SpecialCharTok{+}\NormalTok{ weathersit }\SpecialCharTok{+} 
\NormalTok{              holiday }\SpecialCharTok{+}\NormalTok{ workingday, }
            \AttributeTok{data =}\NormalTok{ data\_model)}

\CommentTok{\# Model summary}
\FunctionTok{summary}\NormalTok{(model)}
\end{Highlighting}
\end{Shaded}

\begin{verbatim}
## 
## Call:
## lm(formula = cnt ~ temp + hum + windspeed + season + yr + weathersit + 
##     holiday + workingday, data = data_model)
## 
## Residuals:
##     Min      1Q  Median      3Q     Max 
## -3677.3  -378.9    67.4   475.7  3347.1 
## 
## Coefficients:
##                     Estimate Std. Error t value Pr(>|t|)    
## (Intercept)          1586.95     226.75   6.999 5.93e-12 ***
## temp                 5108.07     306.98  16.640  < 2e-16 ***
## hum                 -1325.64     295.07  -4.493 8.19e-06 ***
## windspeed           -2795.28     428.06  -6.530 1.24e-10 ***
## seasonSummer         1137.94     113.55  10.021  < 2e-16 ***
## seasonFall            844.10     149.97   5.628 2.61e-08 ***
## seasonWinter         1545.59      96.67  15.988  < 2e-16 ***
## yr2012               2013.71      61.75  32.610  < 2e-16 ***
## weathersitMist       -426.72      81.24  -5.253 1.98e-07 ***
## weathersitLightRain -1916.95     207.31  -9.247  < 2e-16 ***
## holiday1             -624.01     188.44  -3.311 0.000974 ***
## workingday1           118.19      67.92   1.740 0.082232 .  
## ---
## Signif. codes:  0 '***' 0.001 '**' 0.01 '*' 0.05 '.' 0.1 ' ' 1
## 
## Residual standard error: 822.2 on 719 degrees of freedom
## Multiple R-squared:  0.8226, Adjusted R-squared:  0.8199 
## F-statistic: 303.1 on 11 and 719 DF,  p-value: < 2.2e-16
\end{verbatim}

\begin{Shaded}
\begin{Highlighting}[]
\CommentTok{\# Extract key information}
\FunctionTok{cat}\NormalTok{(}\StringTok{"}\SpecialCharTok{\textbackslash{}n}\StringTok{=== INITIAL INTERPRETATION ===}\SpecialCharTok{\textbackslash{}n}\StringTok{"}\NormalTok{)}
\end{Highlighting}
\end{Shaded}

\begin{verbatim}
## 
## === INITIAL INTERPRETATION ===
\end{verbatim}

\begin{Shaded}
\begin{Highlighting}[]
\FunctionTok{cat}\NormalTok{(}\StringTok{"R{-}squared:"}\NormalTok{, }\FunctionTok{round}\NormalTok{(}\FunctionTok{summary}\NormalTok{(model)}\SpecialCharTok{$}\NormalTok{r.squared, }\DecValTok{4}\NormalTok{), }\StringTok{"}\SpecialCharTok{\textbackslash{}n}\StringTok{"}\NormalTok{)}
\end{Highlighting}
\end{Shaded}

\begin{verbatim}
## R-squared: 0.8226
\end{verbatim}

\begin{Shaded}
\begin{Highlighting}[]
\FunctionTok{cat}\NormalTok{(}\StringTok{"Adjusted R{-}squared:"}\NormalTok{, }\FunctionTok{round}\NormalTok{(}\FunctionTok{summary}\NormalTok{(model)}\SpecialCharTok{$}\NormalTok{adj.r.squared, }\DecValTok{4}\NormalTok{), }\StringTok{"}\SpecialCharTok{\textbackslash{}n}\StringTok{"}\NormalTok{)}
\end{Highlighting}
\end{Shaded}

\begin{verbatim}
## Adjusted R-squared: 0.8199
\end{verbatim}

\begin{Shaded}
\begin{Highlighting}[]
\FunctionTok{cat}\NormalTok{(}\StringTok{"F{-}statistic:"}\NormalTok{, }\FunctionTok{round}\NormalTok{(}\FunctionTok{summary}\NormalTok{(model)}\SpecialCharTok{$}\NormalTok{fstatistic[}\DecValTok{1}\NormalTok{], }\DecValTok{4}\NormalTok{), }\StringTok{"}\SpecialCharTok{\textbackslash{}n}\StringTok{"}\NormalTok{)}
\end{Highlighting}
\end{Shaded}

\begin{verbatim}
## F-statistic: 303.0791
\end{verbatim}

\begin{Shaded}
\begin{Highlighting}[]
\CommentTok{\# MEMORY }\AlertTok{NOTE}\CommentTok{:}
\CommentTok{\# {-} R{-}squared indicates what percentage of variance we explain}
\CommentTok{\# {-} Adjusted R{-}squared penalizes for number of variables}
\CommentTok{\# {-} F{-}statistic tests if the model is globally significant (p{-}value \textless{} 0.05)}
\end{Highlighting}
\end{Shaded}

\subsection{Interpretation of
coefficients}\label{interpretation-of-coefficients}

\begin{Shaded}
\begin{Highlighting}[]
\CommentTok{\# Extract coefficients}
\NormalTok{coef\_table }\OtherTok{\textless{}{-}} \FunctionTok{as.data.frame}\NormalTok{(}\FunctionTok{round}\NormalTok{(}\FunctionTok{coef}\NormalTok{(}\FunctionTok{summary}\NormalTok{(model)), }\DecValTok{4}\NormalTok{))}
\NormalTok{coef\_table}\SpecialCharTok{$}\NormalTok{Variable }\OtherTok{\textless{}{-}} \FunctionTok{rownames}\NormalTok{(coef\_table)}
\FunctionTok{rownames}\NormalTok{(coef\_table) }\OtherTok{\textless{}{-}} \ConstantTok{NULL}
\NormalTok{coef\_table }\OtherTok{\textless{}{-}}\NormalTok{ coef\_table[, }\FunctionTok{c}\NormalTok{(}\DecValTok{5}\NormalTok{, }\DecValTok{1}\NormalTok{, }\DecValTok{2}\NormalTok{, }\DecValTok{3}\NormalTok{, }\DecValTok{4}\NormalTok{)]}
\FunctionTok{colnames}\NormalTok{(coef\_table) }\OtherTok{\textless{}{-}} \FunctionTok{c}\NormalTok{(}\StringTok{"Variable"}\NormalTok{, }\StringTok{"Coefficient"}\NormalTok{, }\StringTok{"Std.Error"}\NormalTok{, }\StringTok{"t{-}value"}\NormalTok{, }\StringTok{"p{-}value"}\NormalTok{)}

\FunctionTok{print}\NormalTok{(coef\_table)}
\end{Highlighting}
\end{Shaded}

\begin{verbatim}
##               Variable Coefficient Std.Error t-value p-value
## 1          (Intercept)   1586.9452  226.7473  6.9987  0.0000
## 2                 temp   5108.0686  306.9827 16.6396  0.0000
## 3                  hum  -1325.6366  295.0664 -4.4927  0.0000
## 4            windspeed  -2795.2845  428.0615 -6.5301  0.0000
## 5         seasonSummer   1137.9377  113.5537 10.0211  0.0000
## 6           seasonFall    844.1024  149.9701  5.6285  0.0000
## 7         seasonWinter   1545.5895   96.6724 15.9879  0.0000
## 8               yr2012   2013.7111   61.7519 32.6097  0.0000
## 9       weathersitMist   -426.7212   81.2383 -5.2527  0.0000
## 10 weathersitLightRain  -1916.9498  207.3096 -9.2468  0.0000
## 11            holiday1   -624.0107  188.4392 -3.3115  0.0010
## 12         workingday1    118.1934   67.9151  1.7403  0.0822
\end{verbatim}

\begin{Shaded}
\begin{Highlighting}[]
\FunctionTok{cat}\NormalTok{(}\StringTok{"}\SpecialCharTok{\textbackslash{}n}\StringTok{=== INTERPRETATION OF COEFFICIENTS ===}\SpecialCharTok{\textbackslash{}n\textbackslash{}n}\StringTok{"}\NormalTok{)}
\end{Highlighting}
\end{Shaded}

\begin{verbatim}
## 
## === INTERPRETATION OF COEFFICIENTS ===
\end{verbatim}

\begin{Shaded}
\begin{Highlighting}[]
\FunctionTok{cat}\NormalTok{(}\StringTok{"Intercept (β₀):}\SpecialCharTok{\textbackslash{}n}\StringTok{"}\NormalTok{)}
\end{Highlighting}
\end{Shaded}

\begin{verbatim}
## Intercept (β₀):
\end{verbatim}

\begin{Shaded}
\begin{Highlighting}[]
\FunctionTok{cat}\NormalTok{(}\StringTok{"  Value: "}\NormalTok{, }\FunctionTok{round}\NormalTok{(}\FunctionTok{coef}\NormalTok{(model)[}\DecValTok{1}\NormalTok{], }\DecValTok{2}\NormalTok{), }\StringTok{"}\SpecialCharTok{\textbackslash{}n}\StringTok{"}\NormalTok{)}
\end{Highlighting}
\end{Shaded}

\begin{verbatim}
##   Value:  1586.95
\end{verbatim}

\begin{Shaded}
\begin{Highlighting}[]
\FunctionTok{cat}\NormalTok{(}\StringTok{"  Meaning: Expected number of rentals when all continuous variables are 0}\SpecialCharTok{\textbackslash{}n}\StringTok{"}\NormalTok{)}
\end{Highlighting}
\end{Shaded}

\begin{verbatim}
##   Meaning: Expected number of rentals when all continuous variables are 0
\end{verbatim}

\begin{Shaded}
\begin{Highlighting}[]
\FunctionTok{cat}\NormalTok{(}\StringTok{"  (Has no practical interpretation in this case)}\SpecialCharTok{\textbackslash{}n\textbackslash{}n}\StringTok{"}\NormalTok{)}
\end{Highlighting}
\end{Shaded}

\begin{verbatim}
##   (Has no practical interpretation in this case)
\end{verbatim}

\begin{Shaded}
\begin{Highlighting}[]
\FunctionTok{cat}\NormalTok{(}\StringTok{"Temperature (temp):}\SpecialCharTok{\textbackslash{}n}\StringTok{"}\NormalTok{)}
\end{Highlighting}
\end{Shaded}

\begin{verbatim}
## Temperature (temp):
\end{verbatim}

\begin{Shaded}
\begin{Highlighting}[]
\FunctionTok{cat}\NormalTok{(}\StringTok{"  Coefficient: "}\NormalTok{, }\FunctionTok{round}\NormalTok{(}\FunctionTok{coef}\NormalTok{(model)[}\StringTok{"temp"}\NormalTok{], }\DecValTok{2}\NormalTok{), }\StringTok{"}\SpecialCharTok{\textbackslash{}n}\StringTok{"}\NormalTok{)}
\end{Highlighting}
\end{Shaded}

\begin{verbatim}
##   Coefficient:  5108.07
\end{verbatim}

\begin{Shaded}
\begin{Highlighting}[]
\FunctionTok{cat}\NormalTok{(}\StringTok{"  Interpretation: For each increase of 1 unit in normalized temperature,}\SpecialCharTok{\textbackslash{}n}\StringTok{"}\NormalTok{)}
\end{Highlighting}
\end{Shaded}

\begin{verbatim}
##   Interpretation: For each increase of 1 unit in normalized temperature,
\end{verbatim}

\begin{Shaded}
\begin{Highlighting}[]
\FunctionTok{cat}\NormalTok{(}\StringTok{"                  the number of rentals increases by approx."}\NormalTok{, }\FunctionTok{round}\NormalTok{(}\FunctionTok{coef}\NormalTok{(model)[}\StringTok{"temp"}\NormalTok{], }\DecValTok{0}\NormalTok{), }\StringTok{"bicycles}\SpecialCharTok{\textbackslash{}n}\StringTok{"}\NormalTok{)}
\end{Highlighting}
\end{Shaded}

\begin{verbatim}
##                   the number of rentals increases by approx. 5108 bicycles
\end{verbatim}

\begin{Shaded}
\begin{Highlighting}[]
\FunctionTok{cat}\NormalTok{(}\StringTok{"  Significance: "}\NormalTok{, }\ControlFlowTok{if}\NormalTok{(}\FunctionTok{summary}\NormalTok{(model)}\SpecialCharTok{$}\NormalTok{coefficients[}\StringTok{"temp"}\NormalTok{, }\StringTok{"Pr(\textgreater{}|t|)"}\NormalTok{] }\SpecialCharTok{\textless{}} \FloatTok{0.05}\NormalTok{) }
                            \StringTok{"YES (p \textless{} 0.05)"} \ControlFlowTok{else} \StringTok{"NO"}\NormalTok{, }\StringTok{"}\SpecialCharTok{\textbackslash{}n\textbackslash{}n}\StringTok{"}\NormalTok{)}
\end{Highlighting}
\end{Shaded}

\begin{verbatim}
##   Significance:  YES (p < 0.05)
\end{verbatim}

\begin{Shaded}
\begin{Highlighting}[]
\FunctionTok{cat}\NormalTok{(}\StringTok{"Humidity (hum):}\SpecialCharTok{\textbackslash{}n}\StringTok{"}\NormalTok{)}
\end{Highlighting}
\end{Shaded}

\begin{verbatim}
## Humidity (hum):
\end{verbatim}

\begin{Shaded}
\begin{Highlighting}[]
\FunctionTok{cat}\NormalTok{(}\StringTok{"  Coefficient: "}\NormalTok{, }\FunctionTok{round}\NormalTok{(}\FunctionTok{coef}\NormalTok{(model)[}\StringTok{"hum"}\NormalTok{], }\DecValTok{2}\NormalTok{), }\StringTok{"}\SpecialCharTok{\textbackslash{}n}\StringTok{"}\NormalTok{)}
\end{Highlighting}
\end{Shaded}

\begin{verbatim}
##   Coefficient:  -1325.64
\end{verbatim}

\begin{Shaded}
\begin{Highlighting}[]
\FunctionTok{cat}\NormalTok{(}\StringTok{"  Interpretation: For each increase of 1 unit in humidity,}\SpecialCharTok{\textbackslash{}n}\StringTok{"}\NormalTok{)}
\end{Highlighting}
\end{Shaded}

\begin{verbatim}
##   Interpretation: For each increase of 1 unit in humidity,
\end{verbatim}

\begin{Shaded}
\begin{Highlighting}[]
\FunctionTok{cat}\NormalTok{(}\StringTok{"                  the number of rentals decreases by approx."}\NormalTok{, }\FunctionTok{abs}\NormalTok{(}\FunctionTok{round}\NormalTok{(}\FunctionTok{coef}\NormalTok{(model)[}\StringTok{"hum"}\NormalTok{], }\DecValTok{0}\NormalTok{)), }\StringTok{"bicycles}\SpecialCharTok{\textbackslash{}n}\StringTok{"}\NormalTok{)}
\end{Highlighting}
\end{Shaded}

\begin{verbatim}
##                   the number of rentals decreases by approx. 1326 bicycles
\end{verbatim}

\begin{Shaded}
\begin{Highlighting}[]
\FunctionTok{cat}\NormalTok{(}\StringTok{"  Significance: "}\NormalTok{, }\ControlFlowTok{if}\NormalTok{(}\FunctionTok{summary}\NormalTok{(model)}\SpecialCharTok{$}\NormalTok{coefficients[}\StringTok{"hum"}\NormalTok{, }\StringTok{"Pr(\textgreater{}|t|)"}\NormalTok{] }\SpecialCharTok{\textless{}} \FloatTok{0.05}\NormalTok{) }
                            \StringTok{"YES (p \textless{} 0.05)"} \ControlFlowTok{else} \StringTok{"NO"}\NormalTok{, }\StringTok{"}\SpecialCharTok{\textbackslash{}n\textbackslash{}n}\StringTok{"}\NormalTok{)}
\end{Highlighting}
\end{Shaded}

\begin{verbatim}
##   Significance:  YES (p < 0.05)
\end{verbatim}

\begin{Shaded}
\begin{Highlighting}[]
\FunctionTok{cat}\NormalTok{(}\StringTok{"Wind speed (windspeed):}\SpecialCharTok{\textbackslash{}n}\StringTok{"}\NormalTok{)}
\end{Highlighting}
\end{Shaded}

\begin{verbatim}
## Wind speed (windspeed):
\end{verbatim}

\begin{Shaded}
\begin{Highlighting}[]
\FunctionTok{cat}\NormalTok{(}\StringTok{"  Coefficient: "}\NormalTok{, }\FunctionTok{round}\NormalTok{(}\FunctionTok{coef}\NormalTok{(model)[}\StringTok{"windspeed"}\NormalTok{], }\DecValTok{2}\NormalTok{), }\StringTok{"}\SpecialCharTok{\textbackslash{}n}\StringTok{"}\NormalTok{)}
\end{Highlighting}
\end{Shaded}

\begin{verbatim}
##   Coefficient:  -2795.28
\end{verbatim}

\begin{Shaded}
\begin{Highlighting}[]
\FunctionTok{cat}\NormalTok{(}\StringTok{"  Interpretation: For each increase of 1 unit in wind speed,}\SpecialCharTok{\textbackslash{}n}\StringTok{"}\NormalTok{)}
\end{Highlighting}
\end{Shaded}

\begin{verbatim}
##   Interpretation: For each increase of 1 unit in wind speed,
\end{verbatim}

\begin{Shaded}
\begin{Highlighting}[]
\FunctionTok{cat}\NormalTok{(}\StringTok{"                  the number of rentals decreases by approx."}\NormalTok{, }\FunctionTok{abs}\NormalTok{(}\FunctionTok{round}\NormalTok{(}\FunctionTok{coef}\NormalTok{(model)[}\StringTok{"windspeed"}\NormalTok{], }\DecValTok{0}\NormalTok{)), }\StringTok{"bicycles}\SpecialCharTok{\textbackslash{}n}\StringTok{"}\NormalTok{)}
\end{Highlighting}
\end{Shaded}

\begin{verbatim}
##                   the number of rentals decreases by approx. 2795 bicycles
\end{verbatim}

\begin{Shaded}
\begin{Highlighting}[]
\FunctionTok{cat}\NormalTok{(}\StringTok{"  Significance: "}\NormalTok{, }\ControlFlowTok{if}\NormalTok{(}\FunctionTok{summary}\NormalTok{(model)}\SpecialCharTok{$}\NormalTok{coefficients[}\StringTok{"windspeed"}\NormalTok{, }\StringTok{"Pr(\textgreater{}|t|)"}\NormalTok{] }\SpecialCharTok{\textless{}} \FloatTok{0.05}\NormalTok{) }
                            \StringTok{"YES (p \textless{} 0.05)"} \ControlFlowTok{else} \StringTok{"NO"}\NormalTok{, }\StringTok{"}\SpecialCharTok{\textbackslash{}n\textbackslash{}n}\StringTok{"}\NormalTok{)}
\end{Highlighting}
\end{Shaded}

\begin{verbatim}
##   Significance:  YES (p < 0.05)
\end{verbatim}

\begin{Shaded}
\begin{Highlighting}[]
\FunctionTok{cat}\NormalTok{(}\StringTok{"Season (categorical variables):}\SpecialCharTok{\textbackslash{}n}\StringTok{"}\NormalTok{)}
\end{Highlighting}
\end{Shaded}

\begin{verbatim}
## Season (categorical variables):
\end{verbatim}

\begin{Shaded}
\begin{Highlighting}[]
\ControlFlowTok{for}\NormalTok{(season }\ControlFlowTok{in} \FunctionTok{c}\NormalTok{(}\StringTok{"seasonSummer"}\NormalTok{, }\StringTok{"seasonFall"}\NormalTok{, }\StringTok{"seasonWinter"}\NormalTok{)) \{}
  \ControlFlowTok{if}\NormalTok{(}\SpecialCharTok{!}\FunctionTok{is.na}\NormalTok{(}\FunctionTok{coef}\NormalTok{(model)[season])) \{}
    \FunctionTok{cat}\NormalTok{(}\StringTok{"  "}\NormalTok{, season, }\StringTok{": "}\NormalTok{, }\FunctionTok{round}\NormalTok{(}\FunctionTok{coef}\NormalTok{(model)[season], }\DecValTok{0}\NormalTok{), }\StringTok{"bicycles vs. Spring}\SpecialCharTok{\textbackslash{}n}\StringTok{"}\NormalTok{, }\AttributeTok{sep=}\StringTok{""}\NormalTok{)}
\NormalTok{  \}}
\NormalTok{\}}
\end{Highlighting}
\end{Shaded}

\begin{verbatim}
##   seasonSummer: 1138bicycles vs. Spring
##   seasonFall: 844bicycles vs. Spring
##   seasonWinter: 1546bicycles vs. Spring
\end{verbatim}

\begin{Shaded}
\begin{Highlighting}[]
\FunctionTok{cat}\NormalTok{(}\StringTok{"}\SpecialCharTok{\textbackslash{}n}\StringTok{"}\NormalTok{)}
\end{Highlighting}
\end{Shaded}

\begin{Shaded}
\begin{Highlighting}[]
\FunctionTok{cat}\NormalTok{(}\StringTok{"Year (yr2012):}\SpecialCharTok{\textbackslash{}n}\StringTok{"}\NormalTok{)}
\end{Highlighting}
\end{Shaded}

\begin{verbatim}
## Year (yr2012):
\end{verbatim}

\begin{Shaded}
\begin{Highlighting}[]
\FunctionTok{cat}\NormalTok{(}\StringTok{"  Coefficient: "}\NormalTok{, }\FunctionTok{round}\NormalTok{(}\FunctionTok{coef}\NormalTok{(model)[}\StringTok{"yr2012"}\NormalTok{], }\DecValTok{0}\NormalTok{), }\StringTok{"}\SpecialCharTok{\textbackslash{}n}\StringTok{"}\NormalTok{)}
\end{Highlighting}
\end{Shaded}

\begin{verbatim}
##   Coefficient:  2014
\end{verbatim}

\begin{Shaded}
\begin{Highlighting}[]
\FunctionTok{cat}\NormalTok{(}\StringTok{"  Interpretation: In 2012 there are approx."}\NormalTok{, }\FunctionTok{round}\NormalTok{(}\FunctionTok{coef}\NormalTok{(model)[}\StringTok{"yr2012"}\NormalTok{], }\DecValTok{0}\NormalTok{), }
    \StringTok{"more rentals than in 2011}\SpecialCharTok{\textbackslash{}n\textbackslash{}n}\StringTok{"}\NormalTok{)}
\end{Highlighting}
\end{Shaded}

\begin{verbatim}
##   Interpretation: In 2012 there are approx. 2014 more rentals than in 2011
\end{verbatim}

\begin{Shaded}
\begin{Highlighting}[]
\FunctionTok{cat}\NormalTok{(}\StringTok{"Weather condition (weathersit):}\SpecialCharTok{\textbackslash{}n}\StringTok{"}\NormalTok{)}
\end{Highlighting}
\end{Shaded}

\begin{verbatim}
## Weather condition (weathersit):
\end{verbatim}

\begin{Shaded}
\begin{Highlighting}[]
\ControlFlowTok{for}\NormalTok{(weather }\ControlFlowTok{in} \FunctionTok{c}\NormalTok{(}\StringTok{"weathersitMist"}\NormalTok{, }\StringTok{"weathersitLightRain"}\NormalTok{, }\StringTok{"weathersitHeavyRain"}\NormalTok{)) \{}
  \ControlFlowTok{if}\NormalTok{(}\SpecialCharTok{!}\FunctionTok{is.na}\NormalTok{(}\FunctionTok{coef}\NormalTok{(model)[weather])) \{}
    \FunctionTok{cat}\NormalTok{(}\StringTok{"  "}\NormalTok{, weather, }\StringTok{": "}\NormalTok{, }\FunctionTok{round}\NormalTok{(}\FunctionTok{coef}\NormalTok{(model)[weather], }\DecValTok{0}\NormalTok{), }
        \StringTok{"bicycles vs. Clear weather}\SpecialCharTok{\textbackslash{}n}\StringTok{"}\NormalTok{, }\AttributeTok{sep=}\StringTok{""}\NormalTok{)}
\NormalTok{  \}}
\NormalTok{\}}
\end{Highlighting}
\end{Shaded}

\begin{verbatim}
##   weathersitMist: -427bicycles vs. Clear weather
##   weathersitLightRain: -1917bicycles vs. Clear weather
\end{verbatim}

\subsection{Model diagnostics}\label{model-diagnostics}

\begin{Shaded}
\begin{Highlighting}[]
\CommentTok{\# IMPORTANT: Before interpreting the coefficients, we must verify that}
\CommentTok{\# the assumptions of linear regression are met}

\FunctionTok{par}\NormalTok{(}\AttributeTok{mfrow =} \FunctionTok{c}\NormalTok{(}\DecValTok{2}\NormalTok{, }\DecValTok{2}\NormalTok{))}
\FunctionTok{plot}\NormalTok{(model)}
\end{Highlighting}
\end{Shaded}

\pandocbounded{\includegraphics[keepaspectratio]{Practica_4_Interpretability_EN_files/figure-latex/residuals_diagnostics-1.pdf}}

\begin{Shaded}
\begin{Highlighting}[]
\FunctionTok{par}\NormalTok{(}\AttributeTok{mfrow =} \FunctionTok{c}\NormalTok{(}\DecValTok{1}\NormalTok{, }\DecValTok{1}\NormalTok{))}

\CommentTok{\# MEMORY }\AlertTok{NOTE}\CommentTok{:}
\CommentTok{\# We verify:}
\CommentTok{\# 1. Residuals vs Fitted: Are the errors random? (horizontal line)}
\CommentTok{\# 2. Normal Q{-}Q: Are the residuals normal? (points on the line)}
\CommentTok{\# 3. Scale{-}Location: Is there homoscedasticity? (constant variance)}
\CommentTok{\# 4. Residuals vs Leverage: Are there influential outliers?}
\end{Highlighting}
\end{Shaded}

\begin{Shaded}
\begin{Highlighting}[]
\CommentTok{\# Shapiro{-}Wilk test for normality of residuals}
\NormalTok{shapiro\_test }\OtherTok{\textless{}{-}} \FunctionTok{shapiro.test}\NormalTok{(}\FunctionTok{residuals}\NormalTok{(model))}
\FunctionTok{cat}\NormalTok{(}\StringTok{"Shapiro{-}Wilk test for normality of residuals:}\SpecialCharTok{\textbackslash{}n}\StringTok{"}\NormalTok{)}
\end{Highlighting}
\end{Shaded}

\begin{verbatim}
## Shapiro-Wilk test for normality of residuals:
\end{verbatim}

\begin{Shaded}
\begin{Highlighting}[]
\FunctionTok{cat}\NormalTok{(}\StringTok{"p{-}value:"}\NormalTok{, }\FunctionTok{round}\NormalTok{(shapiro\_test}\SpecialCharTok{$}\NormalTok{p.value, }\DecValTok{4}\NormalTok{), }\StringTok{"}\SpecialCharTok{\textbackslash{}n}\StringTok{"}\NormalTok{)}
\end{Highlighting}
\end{Shaded}

\begin{verbatim}
## p-value: 0
\end{verbatim}

\begin{Shaded}
\begin{Highlighting}[]
\ControlFlowTok{if}\NormalTok{(shapiro\_test}\SpecialCharTok{$}\NormalTok{p.value }\SpecialCharTok{\textgreater{}} \FloatTok{0.05}\NormalTok{) \{}
  \FunctionTok{cat}\NormalTok{(}\StringTok{"CONCLUSION: Residuals are approximately normal (p \textgreater{} 0.05)}\SpecialCharTok{\textbackslash{}n}\StringTok{"}\NormalTok{)}
\NormalTok{\} }\ControlFlowTok{else}\NormalTok{ \{}
  \FunctionTok{cat}\NormalTok{(}\StringTok{"CONCLUSION: Residuals are NOT normal (p \textless{} 0.05)}\SpecialCharTok{\textbackslash{}n}\StringTok{"}\NormalTok{)}
\NormalTok{\}}
\end{Highlighting}
\end{Shaded}

\begin{verbatim}
## CONCLUSION: Residuals are NOT normal (p < 0.05)
\end{verbatim}

\begin{Shaded}
\begin{Highlighting}[]
\CommentTok{\# Breusch{-}Pagan test for homoscedasticity}
\NormalTok{bp\_test }\OtherTok{\textless{}{-}} \FunctionTok{bptest}\NormalTok{(model)}
\FunctionTok{cat}\NormalTok{(}\StringTok{"}\SpecialCharTok{\textbackslash{}n}\StringTok{Breusch{-}Pagan test for homoscedasticity:}\SpecialCharTok{\textbackslash{}n}\StringTok{"}\NormalTok{)}
\end{Highlighting}
\end{Shaded}

\begin{verbatim}
## 
## Breusch-Pagan test for homoscedasticity:
\end{verbatim}

\begin{Shaded}
\begin{Highlighting}[]
\FunctionTok{cat}\NormalTok{(}\StringTok{"p{-}value:"}\NormalTok{, }\FunctionTok{round}\NormalTok{(bp\_test}\SpecialCharTok{$}\NormalTok{p.value, }\DecValTok{4}\NormalTok{), }\StringTok{"}\SpecialCharTok{\textbackslash{}n}\StringTok{"}\NormalTok{)}
\end{Highlighting}
\end{Shaded}

\begin{verbatim}
## p-value: 0
\end{verbatim}

\begin{Shaded}
\begin{Highlighting}[]
\ControlFlowTok{if}\NormalTok{(bp\_test}\SpecialCharTok{$}\NormalTok{p.value }\SpecialCharTok{\textgreater{}} \FloatTok{0.05}\NormalTok{) \{}
  \FunctionTok{cat}\NormalTok{(}\StringTok{"CONCLUSION: There is homoscedasticity (p \textgreater{} 0.05)}\SpecialCharTok{\textbackslash{}n}\StringTok{"}\NormalTok{)}
\NormalTok{\} }\ControlFlowTok{else}\NormalTok{ \{}
  \FunctionTok{cat}\NormalTok{(}\StringTok{"CONCLUSION: There is NO homoscedasticity (p \textless{} 0.05)}\SpecialCharTok{\textbackslash{}n}\StringTok{"}\NormalTok{)}
\NormalTok{\}}
\end{Highlighting}
\end{Shaded}

\begin{verbatim}
## CONCLUSION: There is NO homoscedasticity (p < 0.05)
\end{verbatim}

\begin{Shaded}
\begin{Highlighting}[]
\CommentTok{\# VIF (Variance Inflation Factor) to detect multicollinearity}
\NormalTok{vif\_values }\OtherTok{\textless{}{-}}\NormalTok{ car}\SpecialCharTok{::}\FunctionTok{vif}\NormalTok{(model)}
\FunctionTok{cat}\NormalTok{(}\StringTok{"}\SpecialCharTok{\textbackslash{}n}\StringTok{VIF (Variance Inflation Factor) {-} Multicollinearity:}\SpecialCharTok{\textbackslash{}n}\StringTok{"}\NormalTok{)}
\end{Highlighting}
\end{Shaded}

\begin{verbatim}
## 
## VIF (Variance Inflation Factor) - Multicollinearity:
\end{verbatim}

\begin{Shaded}
\begin{Highlighting}[]
\FunctionTok{print}\NormalTok{(}\FunctionTok{round}\NormalTok{(vif\_values, }\DecValTok{3}\NormalTok{))}
\end{Highlighting}
\end{Shaded}

\begin{verbatim}
##             GVIF Df GVIF^(1/(2*Df))
## temp       3.410  1           1.847
## hum        1.907  1           1.381
## windspeed  1.189  1           1.090
## season     3.538  3           1.234
## yr         1.031  1           1.015
## weathersit 1.792  2           1.157
## holiday    1.071  1           1.035
## workingday 1.078  1           1.038
\end{verbatim}

\begin{Shaded}
\begin{Highlighting}[]
\FunctionTok{cat}\NormalTok{(}\StringTok{"}\SpecialCharTok{\textbackslash{}n}\StringTok{Rule: VIF \textgreater{} 5{-}10 indicates multicollinearity problems}\SpecialCharTok{\textbackslash{}n}\StringTok{"}\NormalTok{)}
\end{Highlighting}
\end{Shaded}

\begin{verbatim}
## 
## Rule: VIF > 5-10 indicates multicollinearity problems
\end{verbatim}

\subsection{Visualization of effects}\label{visualization-of-effects}

\begin{Shaded}
\begin{Highlighting}[]
\CommentTok{\# Extract significant coefficients to visualize}
\NormalTok{coef\_df }\OtherTok{\textless{}{-}} \FunctionTok{data.frame}\NormalTok{(}
  \AttributeTok{variable =} \FunctionTok{names}\NormalTok{(}\FunctionTok{coef}\NormalTok{(model))[}\SpecialCharTok{{-}}\DecValTok{1}\NormalTok{],  }\CommentTok{\# Exclude intercept}
  \AttributeTok{coefficient =} \FunctionTok{coef}\NormalTok{(model)[}\SpecialCharTok{{-}}\DecValTok{1}\NormalTok{],}
  \AttributeTok{se =} \FunctionTok{summary}\NormalTok{(model)}\SpecialCharTok{$}\NormalTok{coefficients[}\SpecialCharTok{{-}}\DecValTok{1}\NormalTok{, }\DecValTok{2}\NormalTok{]}
\NormalTok{)}

\CommentTok{\# Filter only significant coefficients}
\NormalTok{coef\_df}\SpecialCharTok{$}\NormalTok{pvalue }\OtherTok{\textless{}{-}} \FunctionTok{summary}\NormalTok{(model)}\SpecialCharTok{$}\NormalTok{coefficients[}\SpecialCharTok{{-}}\DecValTok{1}\NormalTok{, }\DecValTok{4}\NormalTok{]}
\NormalTok{coef\_df\_sig }\OtherTok{\textless{}{-}}\NormalTok{ coef\_df[coef\_df}\SpecialCharTok{$}\NormalTok{pvalue }\SpecialCharTok{\textless{}} \FloatTok{0.05}\NormalTok{, ]}
\NormalTok{coef\_df\_sig }\OtherTok{\textless{}{-}}\NormalTok{ coef\_df\_sig[}\FunctionTok{order}\NormalTok{(coef\_df\_sig}\SpecialCharTok{$}\NormalTok{coefficient), ]}

\CommentTok{\# Plot}
\FunctionTok{ggplot}\NormalTok{(coef\_df\_sig, }\FunctionTok{aes}\NormalTok{(}\AttributeTok{x =} \FunctionTok{reorder}\NormalTok{(variable, coefficient), }\AttributeTok{y =}\NormalTok{ coefficient)) }\SpecialCharTok{+}
  \FunctionTok{geom\_point}\NormalTok{(}\AttributeTok{size =} \DecValTok{3}\NormalTok{, }\AttributeTok{color =} \StringTok{"steelblue"}\NormalTok{) }\SpecialCharTok{+}
  \FunctionTok{geom\_errorbar}\NormalTok{(}\FunctionTok{aes}\NormalTok{(}\AttributeTok{ymin =}\NormalTok{ coefficient }\SpecialCharTok{{-}} \FloatTok{1.96}\SpecialCharTok{*}\NormalTok{se, }
                    \AttributeTok{ymax =}\NormalTok{ coefficient }\SpecialCharTok{+} \FloatTok{1.96}\SpecialCharTok{*}\NormalTok{se), }
                \AttributeTok{width =} \FloatTok{0.2}\NormalTok{, }\AttributeTok{color =} \StringTok{"steelblue"}\NormalTok{) }\SpecialCharTok{+}
  \FunctionTok{geom\_hline}\NormalTok{(}\AttributeTok{yintercept =} \DecValTok{0}\NormalTok{, }\AttributeTok{linetype =} \StringTok{"dashed"}\NormalTok{, }\AttributeTok{color =} \StringTok{"red"}\NormalTok{) }\SpecialCharTok{+}
  \FunctionTok{coord\_flip}\NormalTok{() }\SpecialCharTok{+}
  \FunctionTok{labs}\NormalTok{(}\AttributeTok{title =} \StringTok{"Linear Model Coefficients (only significant, p \textless{} 0.05)"}\NormalTok{,}
       \AttributeTok{x =} \StringTok{"Variables"}\NormalTok{,}
       \AttributeTok{y =} \StringTok{"Coefficient (95\% CI)"}\NormalTok{) }\SpecialCharTok{+}
  \FunctionTok{theme\_minimal}\NormalTok{() }\SpecialCharTok{+}
  \FunctionTok{theme}\NormalTok{(}\AttributeTok{plot.title =} \FunctionTok{element\_text}\NormalTok{(}\AttributeTok{size =} \DecValTok{12}\NormalTok{, }\AttributeTok{face =} \StringTok{"bold"}\NormalTok{))}
\end{Highlighting}
\end{Shaded}

\pandocbounded{\includegraphics[keepaspectratio]{Practica_4_Interpretability_EN_files/figure-latex/coefficients_plot-1.pdf}}

\begin{Shaded}
\begin{Highlighting}[]
\CommentTok{\# MEMORY }\AlertTok{NOTE}\CommentTok{:}
\CommentTok{\# This graph shows:}
\CommentTok{\# {-} Magnitude of the effect of each variable (Y axis)}
\CommentTok{\# {-} Associated uncertainty (error bars = 95\% CI)}
\CommentTok{\# {-} Only statistically significant variables}
\end{Highlighting}
\end{Shaded}

\begin{Shaded}
\begin{Highlighting}[]
\CommentTok{\# Effects of categorical variables}
\CommentTok{\# Temperature}
\NormalTok{temp\_effect }\OtherTok{\textless{}{-}} \FunctionTok{data.frame}\NormalTok{(}
  \AttributeTok{temp =} \FunctionTok{seq}\NormalTok{(}\DecValTok{0}\NormalTok{, }\DecValTok{1}\NormalTok{, }\FloatTok{0.1}\NormalTok{),}
  \AttributeTok{predicted =} \FunctionTok{predict}\NormalTok{(model, }
                      \AttributeTok{newdata =} \FunctionTok{data.frame}\NormalTok{(}
                        \AttributeTok{temp =} \FunctionTok{seq}\NormalTok{(}\DecValTok{0}\NormalTok{, }\DecValTok{1}\NormalTok{, }\FloatTok{0.1}\NormalTok{),}
                        \AttributeTok{hum =} \FunctionTok{mean}\NormalTok{(data\_model}\SpecialCharTok{$}\NormalTok{hum),}
                        \AttributeTok{windspeed =} \FunctionTok{mean}\NormalTok{(data\_model}\SpecialCharTok{$}\NormalTok{windspeed),}
                        \AttributeTok{season =} \StringTok{"Spring"}\NormalTok{,}
                        \AttributeTok{yr =} \StringTok{"2011"}\NormalTok{,}
                        \AttributeTok{weathersit =} \StringTok{"Clear"}\NormalTok{,}
                        \AttributeTok{holiday =} \StringTok{"0"}\NormalTok{,}
                        \AttributeTok{workingday =} \StringTok{"1"}
\NormalTok{                      ))}
\NormalTok{)}

\NormalTok{p1 }\OtherTok{\textless{}{-}} \FunctionTok{ggplot}\NormalTok{(temp\_effect, }\FunctionTok{aes}\NormalTok{(}\AttributeTok{x =}\NormalTok{ temp, }\AttributeTok{y =}\NormalTok{ predicted)) }\SpecialCharTok{+}
  \FunctionTok{geom\_line}\NormalTok{(}\AttributeTok{color =} \StringTok{"steelblue"}\NormalTok{, }\AttributeTok{size =} \DecValTok{1}\NormalTok{) }\SpecialCharTok{+}
  \FunctionTok{geom\_point}\NormalTok{(}\AttributeTok{size =} \DecValTok{3}\NormalTok{, }\AttributeTok{color =} \StringTok{"steelblue"}\NormalTok{) }\SpecialCharTok{+}
  \FunctionTok{labs}\NormalTok{(}\AttributeTok{title =} \StringTok{"Effect of Temperature on Rentals"}\NormalTok{,}
       \AttributeTok{x =} \StringTok{"Normalized temperature"}\NormalTok{,}
       \AttributeTok{y =} \StringTok{"Predicted rentals"}\NormalTok{) }\SpecialCharTok{+}
  \FunctionTok{theme\_minimal}\NormalTok{() }\SpecialCharTok{+}
  \FunctionTok{theme}\NormalTok{(}\AttributeTok{plot.title =} \FunctionTok{element\_text}\NormalTok{(}\AttributeTok{size =} \DecValTok{11}\NormalTok{, }\AttributeTok{face =} \StringTok{"bold"}\NormalTok{))}

\CommentTok{\# By season}
\NormalTok{season\_effect }\OtherTok{\textless{}{-}} \FunctionTok{data.frame}\NormalTok{(}
  \AttributeTok{season =} \FunctionTok{factor}\NormalTok{(}\FunctionTok{c}\NormalTok{(}\StringTok{"Spring"}\NormalTok{, }\StringTok{"Summer"}\NormalTok{, }\StringTok{"Fall"}\NormalTok{, }\StringTok{"Winter"}\NormalTok{),}
                  \AttributeTok{levels =} \FunctionTok{levels}\NormalTok{(data\_model}\SpecialCharTok{$}\NormalTok{season)),}
  \AttributeTok{predicted =} \FunctionTok{predict}\NormalTok{(model,}
                      \AttributeTok{newdata =} \FunctionTok{data.frame}\NormalTok{(}
                        \AttributeTok{season =} \FunctionTok{factor}\NormalTok{(}\FunctionTok{c}\NormalTok{(}\StringTok{"Spring"}\NormalTok{, }\StringTok{"Summer"}\NormalTok{, }\StringTok{"Fall"}\NormalTok{, }\StringTok{"Winter"}\NormalTok{),}
                                      \AttributeTok{levels =} \FunctionTok{levels}\NormalTok{(data\_model}\SpecialCharTok{$}\NormalTok{season)),}
                        \AttributeTok{temp =} \FunctionTok{mean}\NormalTok{(data\_model}\SpecialCharTok{$}\NormalTok{temp),}
                        \AttributeTok{hum =} \FunctionTok{mean}\NormalTok{(data\_model}\SpecialCharTok{$}\NormalTok{hum),}
                        \AttributeTok{windspeed =} \FunctionTok{mean}\NormalTok{(data\_model}\SpecialCharTok{$}\NormalTok{windspeed),}
                        \AttributeTok{yr =} \StringTok{"2011"}\NormalTok{,}
                        \AttributeTok{weathersit =} \StringTok{"Clear"}\NormalTok{,}
                        \AttributeTok{holiday =} \StringTok{"0"}\NormalTok{,}
                        \AttributeTok{workingday =} \StringTok{"1"}
\NormalTok{                      ))}
\NormalTok{)}

\NormalTok{p2 }\OtherTok{\textless{}{-}} \FunctionTok{ggplot}\NormalTok{(season\_effect, }\FunctionTok{aes}\NormalTok{(}\AttributeTok{x =}\NormalTok{ season, }\AttributeTok{y =}\NormalTok{ predicted, }\AttributeTok{fill =}\NormalTok{ season)) }\SpecialCharTok{+}
  \FunctionTok{geom\_bar}\NormalTok{(}\AttributeTok{stat =} \StringTok{"identity"}\NormalTok{, }\AttributeTok{color =} \StringTok{"black"}\NormalTok{) }\SpecialCharTok{+}
  \FunctionTok{scale\_fill\_manual}\NormalTok{(}\AttributeTok{values =} \FunctionTok{c}\NormalTok{(}\StringTok{"Spring"} \OtherTok{=} \StringTok{"\#90EE90"}\NormalTok{, }
                               \StringTok{"Summer"} \OtherTok{=} \StringTok{"\#FFD700"}\NormalTok{,}
                               \StringTok{"Fall"} \OtherTok{=} \StringTok{"\#FF6347"}\NormalTok{,}
                               \StringTok{"Winter"} \OtherTok{=} \StringTok{"\#87CEEB"}\NormalTok{)) }\SpecialCharTok{+}
  \FunctionTok{labs}\NormalTok{(}\AttributeTok{title =} \StringTok{"Effect of Season on Rentals"}\NormalTok{,}
       \AttributeTok{x =} \StringTok{"Season"}\NormalTok{,}
       \AttributeTok{y =} \StringTok{"Predicted rentals"}\NormalTok{) }\SpecialCharTok{+}
  \FunctionTok{theme\_minimal}\NormalTok{() }\SpecialCharTok{+}
  \FunctionTok{theme}\NormalTok{(}\AttributeTok{plot.title =} \FunctionTok{element\_text}\NormalTok{(}\AttributeTok{size =} \DecValTok{11}\NormalTok{, }\AttributeTok{face =} \StringTok{"bold"}\NormalTok{),}
        \AttributeTok{legend.position =} \StringTok{"none"}\NormalTok{)}

\FunctionTok{grid.arrange}\NormalTok{(p1, p2, }\AttributeTok{ncol =} \DecValTok{2}\NormalTok{)}
\end{Highlighting}
\end{Shaded}

\pandocbounded{\includegraphics[keepaspectratio]{Practica_4_Interpretability_EN_files/figure-latex/categorical_variable_effects-1.pdf}}

\begin{Shaded}
\begin{Highlighting}[]
\CommentTok{\# MEMORY }\AlertTok{NOTE}\CommentTok{:}
\CommentTok{\# These graphs show:}
\CommentTok{\# {-} How model predictions vary with independent variables}
\CommentTok{\# {-} The effect is linear (straight line) for continuous variables}
\CommentTok{\# {-} Marked differences between categories for categorical variables}
\end{Highlighting}
\end{Shaded}

\subsection{Model summary}\label{model-summary}

\begin{Shaded}
\begin{Highlighting}[]
\FunctionTok{cat}\NormalTok{(}\StringTok{"=== EXECUTIVE MODEL SUMMARY ===}\SpecialCharTok{\textbackslash{}n\textbackslash{}n}\StringTok{"}\NormalTok{)}
\end{Highlighting}
\end{Shaded}

\begin{verbatim}
## === EXECUTIVE MODEL SUMMARY ===
\end{verbatim}

\begin{Shaded}
\begin{Highlighting}[]
\FunctionTok{cat}\NormalTok{(}\StringTok{"Goodness of fit:}\SpecialCharTok{\textbackslash{}n}\StringTok{"}\NormalTok{)}
\end{Highlighting}
\end{Shaded}

\begin{verbatim}
## Goodness of fit:
\end{verbatim}

\begin{Shaded}
\begin{Highlighting}[]
\FunctionTok{cat}\NormalTok{(}\StringTok{"{-} R{-}squared:"}\NormalTok{, }\FunctionTok{round}\NormalTok{(}\FunctionTok{summary}\NormalTok{(model)}\SpecialCharTok{$}\NormalTok{r.squared }\SpecialCharTok{*} \DecValTok{100}\NormalTok{, }\DecValTok{1}\NormalTok{), }
    \StringTok{"\% of variance explained}\SpecialCharTok{\textbackslash{}n}\StringTok{"}\NormalTok{)}
\end{Highlighting}
\end{Shaded}

\begin{verbatim}
## - R-squared: 82.3 % of variance explained
\end{verbatim}

\begin{Shaded}
\begin{Highlighting}[]
\FunctionTok{cat}\NormalTok{(}\StringTok{"{-} Adjusted R{-}squared:"}\NormalTok{, }\FunctionTok{round}\NormalTok{(}\FunctionTok{summary}\NormalTok{(model)}\SpecialCharTok{$}\NormalTok{adj.r.squared }\SpecialCharTok{*} \DecValTok{100}\NormalTok{, }\DecValTok{1}\NormalTok{), }\StringTok{"\%}\SpecialCharTok{\textbackslash{}n}\StringTok{"}\NormalTok{)}
\end{Highlighting}
\end{Shaded}

\begin{verbatim}
## - Adjusted R-squared: 82 %
\end{verbatim}

\begin{Shaded}
\begin{Highlighting}[]
\FunctionTok{cat}\NormalTok{(}\StringTok{"{-} RMSE:"}\NormalTok{, }\FunctionTok{round}\NormalTok{(}\FunctionTok{sqrt}\NormalTok{(}\FunctionTok{sum}\NormalTok{(}\FunctionTok{residuals}\NormalTok{(model)}\SpecialCharTok{\^{}}\DecValTok{2}\NormalTok{) }\SpecialCharTok{/} \FunctionTok{df.residual}\NormalTok{(model)), }\DecValTok{2}\NormalTok{), }
    \StringTok{"bicycles}\SpecialCharTok{\textbackslash{}n\textbackslash{}n}\StringTok{"}\NormalTok{)}
\end{Highlighting}
\end{Shaded}

\begin{verbatim}
## - RMSE: 822.16 bicycles
\end{verbatim}

\begin{Shaded}
\begin{Highlighting}[]
\FunctionTok{cat}\NormalTok{(}\StringTok{"Validated assumptions:}\SpecialCharTok{\textbackslash{}n}\StringTok{"}\NormalTok{)}
\end{Highlighting}
\end{Shaded}

\begin{verbatim}
## Validated assumptions:
\end{verbatim}

\begin{Shaded}
\begin{Highlighting}[]
\FunctionTok{cat}\NormalTok{(}\StringTok{"{-} Normality of residuals:"}\NormalTok{, }
    \ControlFlowTok{if}\NormalTok{(shapiro\_test}\SpecialCharTok{$}\NormalTok{p.value }\SpecialCharTok{\textgreater{}} \FloatTok{0.05}\NormalTok{) }\StringTok{"✓ YES"} \ControlFlowTok{else} \StringTok{"✗ NO"}\NormalTok{, }\StringTok{"}\SpecialCharTok{\textbackslash{}n}\StringTok{"}\NormalTok{)}
\end{Highlighting}
\end{Shaded}

\begin{verbatim}
## - Normality of residuals: ✗ NO
\end{verbatim}

\begin{Shaded}
\begin{Highlighting}[]
\FunctionTok{cat}\NormalTok{(}\StringTok{"{-} Homoscedasticity:"}\NormalTok{, }
    \ControlFlowTok{if}\NormalTok{(bp\_test}\SpecialCharTok{$}\NormalTok{p.value }\SpecialCharTok{\textgreater{}} \FloatTok{0.05}\NormalTok{) }\StringTok{"✓ YES"} \ControlFlowTok{else} \StringTok{"✗ NO"}\NormalTok{, }\StringTok{"}\SpecialCharTok{\textbackslash{}n}\StringTok{"}\NormalTok{)}
\end{Highlighting}
\end{Shaded}

\begin{verbatim}
## - Homoscedasticity: ✗ NO
\end{verbatim}

\begin{Shaded}
\begin{Highlighting}[]
\FunctionTok{cat}\NormalTok{(}\StringTok{"{-} Multicollinearity:"}\NormalTok{, }
    \ControlFlowTok{if}\NormalTok{(}\FunctionTok{max}\NormalTok{(vif\_values) }\SpecialCharTok{\textless{}} \DecValTok{5}\NormalTok{) }\StringTok{"✓ Acceptable"} \ControlFlowTok{else} \StringTok{"✗ Problems"}\NormalTok{, }\StringTok{"}\SpecialCharTok{\textbackslash{}n\textbackslash{}n}\StringTok{"}\NormalTok{)}
\end{Highlighting}
\end{Shaded}

\begin{verbatim}
## - Multicollinearity: ✓ Acceptable
\end{verbatim}

\begin{Shaded}
\begin{Highlighting}[]
\FunctionTok{cat}\NormalTok{(}\StringTok{"Most important variables (coefficient \textgreater{} 500 in absolute value):}\SpecialCharTok{\textbackslash{}n}\StringTok{"}\NormalTok{)}
\end{Highlighting}
\end{Shaded}

\begin{verbatim}
## Most important variables (coefficient > 500 in absolute value):
\end{verbatim}

\begin{Shaded}
\begin{Highlighting}[]
\NormalTok{important\_coefs }\OtherTok{\textless{}{-}} \FunctionTok{coef}\NormalTok{(model)[}\SpecialCharTok{{-}}\DecValTok{1}\NormalTok{]}
\NormalTok{important\_coefs }\OtherTok{\textless{}{-}}\NormalTok{ important\_coefs[}\FunctionTok{abs}\NormalTok{(important\_coefs) }\SpecialCharTok{\textgreater{}} \DecValTok{500}\NormalTok{]}
\ControlFlowTok{for}\NormalTok{(var }\ControlFlowTok{in} \FunctionTok{names}\NormalTok{(important\_coefs)) \{}
  \FunctionTok{cat}\NormalTok{(}\StringTok{"  {-}"}\NormalTok{, var, }\StringTok{":"}\NormalTok{, }\FunctionTok{round}\NormalTok{(important\_coefs[var], }\DecValTok{0}\NormalTok{), }\StringTok{"}\SpecialCharTok{\textbackslash{}n}\StringTok{"}\NormalTok{)}
\NormalTok{\}}
\end{Highlighting}
\end{Shaded}

\begin{verbatim}
##   - temp : 5108 
##   - hum : -1326 
##   - windspeed : -2795 
##   - seasonSummer : 1138 
##   - seasonFall : 844 
##   - seasonWinter : 1546 
##   - yr2012 : 2014 
##   - weathersitLightRain : -1917 
##   - holiday1 : -624
\end{verbatim}

\begin{Shaded}
\begin{Highlighting}[]
\CommentTok{\# MEMORY }\AlertTok{NOTE}\CommentTok{:}
\CommentTok{\# The model is relatively good but could be improved:}
\CommentTok{\# {-} Considering non{-}linear relationships}
\CommentTok{\# {-} Adding interaction terms}
\CommentTok{\# {-} Using regularization techniques (Ridge, Lasso)}
\end{Highlighting}
\end{Shaded}

\begin{center}\rule{0.5\linewidth}{0.5pt}\end{center}

\subsection{Notes for the next
section}\label{notes-for-the-next-section}

\textbf{For whoever continues with the next section:}

\begin{enumerate}
\def\labelenumi{\arabic{enumi}.}
\tightlist
\item
  Data is prepared and model fitted
\item
  The model is reasonably interpretable (R² ≈ 0.73)
\item
  There are some categorical variables that could be explored further
\item
  Interaction terms can be considered (e.g: temperature × season)
\item
  The following sections should include:

  \begin{itemize}
  \tightlist
  \item
    Model validation on test data
  \item
    Analysis of relative importance of variables
  \item
    Comparison with alternative models
  \item
    Application to new predictions
  \end{itemize}
\end{enumerate}

\textbf{Code to load this model in the next section:}

\begin{Shaded}
\begin{Highlighting}[]
\CommentTok{\# load("linear\_model\_bike.RData")  \# If the model is saved}
\CommentTok{\# The model is called: model}
\end{Highlighting}
\end{Shaded}


\end{document}
